\section{模型评价与结论}
\subsection{模型特点}
四问以一致的变量和环层通量结构处理常热物性、双向耦合、全域阈值及径向收缩。有限体积保留变系数散度与收支关系，表面加密兼顾薄层分辨和题定输出；解析解、制造解、退化及二维对照分别检查实现与简化。达标时区分临界根、严格不等式与四位显示，收缩时以材料坐标保持干物质和水质量守恒，并按当前物理距离输出结果。

\subsection{局限与可推广方向}
第三问二维校核为指定情景下的径向简化提供了定量依据。第四问结合径向收缩时长径比增大、端侧面积比减小的几何特征沿用这一近似；其依据仍受各向同性及端面条件假设限制，收缩情景下的端面时间误差尚未单独量化。

等效水分边界缺少独立吸附等温关系，热方程也未显式描述汽化潜热和水分携焓。更完整的界面模型需要药材平衡含水率、相变和气固传质数据。仅凭当前方程的收支残差，不能证明这些机制对真实预测可忽略。

经验密度与几何干物质密度分开使用，是本模型的一项有效近似。以问题四为例，若同时把$\rho(U)$视为真实湿密度，则初始干密度为$(760+90\times2.55)/(1+2.55)\approx278.73$ kg/m$^3$。72 h半径对应的体积比为$(1.198/2)^2=0.358801$，固定长度且干物质守恒要求平均干密度约776.84 kg/m$^3$；但对任意$U\ge0$，
\[
\frac{760+90U}{1+U}=90+\frac{670}{1+U}\le760\ \mathrm{kg/m^3}.
\]
这说明上述解释不能同时精确成立，因此本文只将$\rho c_p$用于有效热容量。该检查针对72 h数据和指定解释，不是51.0877 h报告时刻的误差估计，也不说明题目无法建模。

后期环境为恒定延续假设，尚无内部实测数据验证预测。问题三温度情景使时长变化约6\%，后期水分量$\pm10\%$使其变化约0.2\%；这些输入不确定性意味着数值时长末位不能解释为实际预测精度。

\subsection{结论}
在题给初值、经验物性和本文明确的等效假设下，问题一1800 s时中心温度33.5758 $^\circ$C、表面36.7879 $^\circ$C，水分损失主要集中在外层；问题二3 h时中心温度49.8494 $^\circ$C、含水率1.7662 kg/kg，温度接近平衡早于内部达到干燥标准。全域阈值计算给出附录3固定半径模型的报告时长57.4731 h、附录4给定收缩模型的报告时长51.0877 h。二者均经未舍入末态和时空加密核验，且移动边界计算未超出半径观测范围。补充的附录4固定半径对照报告时长为129.8448 h，在相同物性与环境条件下，收缩使模型预测时长缩短78.7571 h，约占固定组的60.65\%。
